\chapter{Over deze vertaling}
\label{cha:translation-about}

\section{Voorwoord bij de vertaling}

De Mega65 is een klassieke homecomputer zoals ze waren in het gouden decennium van de homecomputer, de jaren '80.
Weliswaar heeft de Mega65 veel krachtiger hardware dan de computers van destijds, de filosofie achter de computer probeert
de filosofie achter de klassieke homecomputers zo getrouw mogelijk na te leven. Eén van de elementen waaruit dat blijkt
is de handleiding. Computers uit de jaren '80 werden geleverd met een informatieve handleiding. Een handleiding die begon
bij het prille begin, maar de lezer niet voor dom hield, maar hem opleidde tot een volwaardige programmeur.

De handleiding van de Mega65 doet precies dat: De handleiding begint bij het begin, legt uit hoe de computer gebruikt moet
worden, en begint daarna met de beginselen van het programmeren in BASIC. Er is evenwel nog een niet on-belangrijk detail: De
computers van de jaren '80 hadden niet alleen een informatieve handleiding, zij hadden een informatieve handleiding in
alle talen. Dat deed ertoe. Een tiener die in jaren '80 zijn eerste computer kreeg, was vaak onkundig om een volwaardige
handleiding in het Engels te kunnen lezen. Juist doordat iedere bezitter van een homecomputer een handleiding had in zijn
eigen taal, leerden ontzettend veel van hen te programmeren. De homecomputers uit onze jonge jaren brachten ons niet
uitsluitend amusement, zij hebben ons gevormd tot kundige vakmensen, waardoor we de wereld konden bouwen zoals deze nu is.

De Mega65 is niet uitsluitend bedoeld voor degene die de jaren '80 actief heeft meegemaakt. De Mega65 is ook niet uitsluitend
bedoeld voor computerexperts, nee, de Mega65 wil de fascinerende tijdsgeest van computers in de jaren '80 aan iedereen laten
zien die in computers geïnteresseerd is. Dat maakt de taal van de handleiding een relevant onderwerp.

Het Engels is een taal met een enorme bibliotheek aan kennis en omdat de grammatica relatief eenvoudig is, spreken veel
mensen het. Dat maakt het een logische taal om documentatie in te publiceren. Toch heeft deze taal ook nadelen. Juist
doordat het Engels een eenvoudige grammatica heeft, bevatten zinnen minder informatie en kunnen begrippen zelfs dubbelzinnig
opgevat worden. Dat maakt vertaling vanuit het Engels uitdagend: Wordt een tekst te letterlijk vertaald, dan ontstaat een
onnatuurlijk "vertaalnederlands" dat heel anders is dan het Nederlands dat we dagelijks spreken. Vertaling vanuit het Engels
is een heel creatief proces, waarbij de vertaler zaken als verkleinwoorden en positiewerkwoorden creatief moet toevoegen.

Er bestaat het risico dat de vertaler een dubbelzinnig begrip verkeerd interpreteert en daardoor de lezer wordt misleid.
Een mooi voorbeeld is het begrip "serial pherpheral bus", dat door automatische vertaalhulpmiddelen aanvankelijk vertaald
werd als"seriële perife bus". Deze problemen ontstaan doordat de Engelse grammatica geen onderscheid kan maken tussen een
bijvoeglijk naamwoord en een zelfstandig naamwoord, waardoor de vertaler de mist in kan gaan. Als dit soort fouten niet
worden opgemerkt, moet de lezer in dit soort situaties terugredeneren naar het Engels en concluderen dat "seriële
periferierbus" bedoeld werd, of in nog beter Nederlands, "seriële randapparatuurbus". Er ligt dus een belangrijke
verantwoordelijkheid op de vertaler om goed werk af te leveren, anders mist de vertaling zijn doel.

Toch is dit wellicht het beste argument dat pleit om technische documenten te vertalen: Immers, als een vertaler dit soort
interpretatiefouten kan maken, kan de lezer die fouten ook maken. De lezer zal in praktijk, als hij de materie beter begint
te begrijpen, uiteindelijk vaak alsnog op de juiste betekenis van het woord komen. Desondanks maakt dit het Nederlands een
zeer fijne taal om technsiche materie heel precies in uit te leggen. Het Nederlands is, bijvoorbeeld dankzij de krachtige
woordconstuctiemechanismen, een taal waarin zinnen iets met heel hoge precisie kunnen formuleren, waardoor minder interpretatie
aan de lezer overgelaten wordt. De tekst wordt daardoor toegankelijker en makkelijker begrijpbaar.

Met die insteek heb ik deze vertaling gemaakt. Deze vertaling is geen vertaling die woord-voor-woord het origineel volgt.
In veel gevallen is dat niet eens zin-voor-zin. Bij het vertalen stond de boodschap van een tekstfragment te allen tijde
voorop en de boodschap is opnieuw geformuleerd zoals dit in het Nederlands uitgelegd zou worden, creatief gebruik makend van
de mogelijkheden die onze taal heeft.

Ik hoop dat met deze handleiding u nog meer plezierige en interessante momenten met uw Mega65 zult hebben.
\\~\\
-- {\bf Daniël Mantione}, oktober 2026.

\section{Praktische zaken}

Een bezitter van een MEGA65 is het volste respect waard, dus we zullen u in dit boek met "u" aanspreken.

Deze vertaling is geschreven volgens de Witte Spelling, waarbij de klassieke tussen-n-regel "n = meervoud" gebruikt wordt.
Bij woorden waarbij zowel een n als een s voor het meervoud gebruikt kan worden, verkiezen we de n.

De Mega65-computer zelf spreekt geen Nederlands en gebruikt de Engelse getalnotatie met de punt (.) als decimaal
scheidingsteken. In de handleiding halen we regelmatig fragmenten aan van programma's of andere tekst die ingetypt
wordt aan. Het zou zeer verwarrend zijn als we voortdurend tussen de punt en komma als decimaal
scheidingsteken zouden wisselen, daarom gebruiken we, in afwijking van de Nederlandse spelling, de punt als decimaal
scheidingsteken. Om verwarring te voorkomen gebruiken we geen scheidingsteken voor duizentallen.

Voor de verschillende soorten getallen gebruiken we de termen zoals u ze in de wiskundeles geleerd heeft:
\begin{enumerate}
\item Het begrip {\em unsigned integer} vertalen we met {\bf natuurlijk getal}.
\item Het begrip {\em signed integer} vertalen we met {\bf geheel getal}.
\item Het begrip {\em real number} vertalen we met {\bf reëel getal}.
\end{enumerate}

In delen van de documentatie die toegespitst zijn op beginners zullen we {\em integer} aanvankelijk vermijden
en ook zullen we een {\em string} een {\bf tekenreeks} noemen. In delen van de documentatie die toegespitst zijn op
meer ervaren programmeurs hebben we wat minder moeite met deze termen, ze zijn immers goed ingeburgerd in het
informaticajargon.

\section{Hoe deze vertaling gemaakt is}

Het omvattende boek "Alles over de Mega65" beslaat inmiddels meer dan 1500 bladzijden. Laten we er geen doekjes om winden: Dat
is niet allemaal met de hand vertaald. Deze handleiding is gemaakt met een combinatie van zowel handmatige, programmatische,
als automatische vertaling. Delen met van de handleiding met veel afbeeldingen en codevoorbeelden zijn met de hand vertaald,
waaronder de afbeeldingen en codevoorbeelden zelf.

Voor de naslaghoofdstukken, zoals bijvoorbeeld de BASIC65-naslag en registertabellen, geldt dat veel termen voor iedere
behandeld commando terugkomen. Sommige zinnen komen ettelijke malen terug. Hiervoor zijn programma's geschreven, die de
vertaalde hoofdstukken vanuit tabellen opbouwen, zodat elke term in het gehele boek consistent wordt gehanteerd.

Voor hoofdstukken met grote hoeveelheden "droge tekst" is gebruikt gemaakt van automatische vertaling door middel van
kunstmatige intelligentie. Daartoe is eerst een stijlboek opgesteld en daarna is de tekst vertaald door een taalmodel. Daarna
is met een tweede taalmodel gecontroleerd of de tekst aan kwaliteitseisen voldoet op het gebied van stijlboek, spelling en of
de tekst als natuurlijk Nederlands leest. Alle automatisch vertaalde tekst is nog eens met de hand nagelopen, gecorrigeerd
en geherformuleerd.
